\begin{itemize}[leftmargin=5mm]
  \setlength{\itemsep}{0pt}
  \item 大量读入浮点数时，即使关闭同步流，\texttt{cin} 仍可能很慢。优先使用
    \texttt{scanf}；若题面数据实际均为整数，则先按整数读入，再转换为
    \texttt{long double} 参与计算。
  \item \texttt{eps} 并非越小越好。输入浮点数本身的精度可能有限；若
    \texttt{eps} 小于输入误差的量级，比较标准会与有效输入精度不匹配，反而导致错误。
  \item 看到 \texttt{acos} 的参数接近 $1$、大数减大数，或
    \texttt{x - sin(x)}、\texttt{x - cos(x)} 这类小量相减时，立即警觉消去误差与数值稳定性。
  \item 处理方向夹角时，可先按极角排序，再检查环上相邻极角之差（包括首尾跨
    $2\pi$ 的一对）。若差值为 $d$，取 $\min(d, 2\pi-d)$；等价地，$d>\pi$
    时改取 $2\pi-d$。
  \item 只要数据与运算可以保持整数形式，就优先使用 \texttt{long long}，不要过早转为浮点数。
  \item 只判闭半平面交非空时使用 \texttt{hp\_feasible(vector<HP>)}；
    求交多边形顶点时使用原半平面交，注意其无界和零面积限制。
    已知 $aX+bY\ge c$ 时直接构造 \texttt{HP(\{a,b\},c)}，构造时统一单位化法向量。
    有向线左侧可用 \texttt{HP(line)}，但两点相减可能损失原本精确的平行关系。
    核心仅补偿浮点舍入，不主动放宽边界；若题目需要距离容差，可对非零法向量的
    \texttt{HP} 手动令 \texttt{b -= tolerance}。舍入补偿不能代替精确算术。
  \item 含 \texttt{eps} 的比较器用于 \texttt{std::sort} 时需特别谨慎：
    “近似相等”可能不具传递性，破坏 strict weak ordering（严格弱序）。
    能用整数或精确方向排序时，优先使用精确比较，并注意中间运算溢出。
  \item 若答案必为整数且可能超过 $10^{19}$，使用 \texttt{\_\_int128} 计算并输出，
    不要使用 \texttt{long double} 代替整数类型；浮点数无法保证大整数运算的精确性。
\end{itemize}
